\documentclass[11pt]{article}
\usepackage[margin=1in]{geometry}
\usepackage{amsmath,amssymb,amsthm,booktabs,longtable}
\usepackage[T1]{fontenc}
\usepackage{hyperref}
\newtheorem{theorem}{Theorem}
\newtheorem{lemma}{Lemma}
\newtheorem{corollary}{Corollary}
\newcommand{\Tr}{\operatorname{Tr}}
\newcommand{\diam}{\operatorname{diam}}
\newcommand{\Norm}{\operatorname{Norm}}
\newcommand{\SWAP}{\operatorname{SWAP}}
\newcommand{\TV}{\operatorname{TV}}
\title{Reusable autonomous processing under spatial Pareto bursts\\
and the optimal asymptotic memory-diameter exponent}
\author{Constructive continuation of the accepted finite-seed programme}
\date{5 November 2026}
\begin{document}
\maketitle

\begin{abstract}
The accepted finite-seed compiler and one-pass autonomous conveyor do not by
themselves justify reusing their hardware: repeatedly decoding the retained
state would incur an interface error on every lap, and an unprotected
one-shot latch would reintroduce a reliable controller. We construct a
reusable processor with one sacrificial input/output lane. A protected
loaded flag admits the once-supplied unknown input on the first lap; a
protected countdown releases it on the selected final lap. The retained
state stays encoded throughout. Arbitrary failures of intermediate ascents
and descents are harmless, by strong recursive normalization and the
identity branch of the encoded conditional swap.

The operation preserves an arbitrary external reference and supports a
repeated finite computation kernel with mutable program and work records.
It uses a charged, noisy, input-independent product preparation for its
classical counter and flags, without an initially encoded unknown quantum
state. At fixed sufficiently small positive intensity, the construction
also yields autonomous quantum memory with full-device diameter
$T^{1/(\alpha-d)+o(1)}$, hardware
$T^{2/(\alpha-d)+o(1)}$, and all-hardware update work
$T^{1+2/(\alpha-d)+o(1)}$. A complete-device erasure proves the diameter
exponent optimal. Hardware optimality and an exact constant-factor
$T^{1/(\alpha-d)}$ diameter bound are not claimed.
\end{abstract}

\section{Question, model, and earned consequence}

The selected question is whether a finite unknown input can undergo a long
computation in a \emph{reused} autonomous local processor under the accepted
correlated spatial faults, including its mutable control, without paying
for an entire new physical circuit layer for every execution step. The
construction below answers this within the same spatially patterned,
synchronous local-channel scope as the accepted autonomous conveyor.

Let $d\geq2$. Burst centers and times form a homogeneous marked Poisson
process on $\mathbb R^d\times\mathbb R$ with intensity $\lambda\,dz\,dt$.
The radius law is
\[
 \Pr(R\geq u)=u^{-\alpha},\qquad u\geq1,
 \qquad r:=\alpha-d>1.
\]
A mark permits an arbitrary joint CPTP corruption on the mutable registers
in its spatial cube. Fault values may be adversarial, share an environment,
and have memory. Geometry remains the declared Poisson source process.
Every data, program, message, countdown, latch, syndrome, and work record
is mutable and exposed. The finite register dimension and per-cell update
rate are bounded independently of the requested horizon.

The following model premises remain explicit.
\begin{enumerate}
\item A fixed finite alphabet of bounded-local channels, including reset,
dephasing, classical reversible computation, and a quantum universal
gate set, is available. Its local laws and static wiring are immutable.
\item A finite input-independent spatial layout may depend on the
requested horizon. Repeated evolution applies one and the same fixed
constant-depth local update to that layout. The proof retains synchronous
microscopic updates; it does not establish an asynchronous or Hamiltonian
implementation.
\item A bounded-depth parallel \emph{noisy} product preparation is available
before the repeated evolution. It sets the canonical physical leaves of
the known classical controller words. Its reset locations and all faults
during its fixed-duration window are charged. The remaining quantum work
registers may start in arbitrary states.
\item The unknown quantum/classical input occupies a fixed finite compact
seed at time zero. Its size is independent of the requested retention
horizon. A one-time local input access and final local output access have
their declared errors $\eta_{\rm in},\eta_{\rm out}$. The device receives
no later input-dependent restoration source.
\end{enumerate}

Known controller initialization is input-independent but can grow with the
horizon. This is a product state preparation of known classical symbols,
not a supplied extensive encoding of the unknown state or answer. Preparing
the device from arbitrary controller junk is a stronger operation and is
not asserted here. Physical fabrication, initialization, execution, and
readout are all included in the resource accounting below.

For the asymptotically exact exponent, the fixed noise parameter is
computable or comes with effective certified bounds, and conservative
constants in the accepted constructive compiler are effectively available.
A fixed certified lower bound on $\alpha$ already gives the corresponding
fixed-family bounds without this refinement.

\begin{theorem}[Reusable processor, fixed padding]\label{thm:reuse}
Fix the finite raw input and output interfaces, the burst law above, a
mixed classical/quantum primitive compiler, and a padding choice
$B$ with $A=CB$ and
\[
 1<\kappa:=\log_B A<r,\qquad
 \sigma:=r-\kappa>0.
\]
For sufficiently small fixed $\lambda>0$, a finite computation kernel of
logical width $w$ and local depth $b$, repeated $K$ times, has a uniform
autonomous realization. Include its $O(\log(K+1))$ countdown bits and all
program, work, and control records in $w$, and put
$\bar b=O(b+w+1)$ and $G=O(w\bar b)$.

At encoding extent $H=B^m$, its reference-inclusive channel error satisfies
\begin{equation}\label{eq:generalerror}
 \frac12\|\mathcal E-\mathcal C\|_\diamond
 \leq \eta_{\rm in}+\eta_{\rm out}
     + C_{\rm io}\lambda
     + C_{\rm init}\lambda w_{\rm c}H^{-r}
     + C_{\rm bulk}\lambda KG H^{-\sigma}
     +\epsilon_{\rm synth}.
\end{equation}
Here $w_{\rm c}=O(\log(K+1))$ counts the known initial controller words,
$\mathcal C$ is the intended $K$-step input/output channel, and all
approximation errors in a chosen logical gate synthesis are included in
$\epsilon_{\rm synth}$. The memory, timer, loading and release logic use
exact finite circuits and do not introduce a per-lap synthesis error.

The complete hardware width, one-period duration, installed register count,
post-input duration, and all-register update work obey
\begin{align}
 W&=O(wH),&
 P&=O(\bar bH^\kappa),&
 N&=O(WP)=O(w\bar bH^{1+\kappa}),\label{eq:res1}\\
 t_{\rm run}&=KP,&
 U_{\rm all}&=O(NKP)
   =O(Kw\bar b^2H^{1+2\kappa}).\label{eq:res2}
\end{align}
The input and output are at the same fixed local port at the first and
selected final period boundary. No mutable controller is exempt from the
noise.
\end{theorem}

Taking
$H=O((1+\lambda KG/\epsilon_{\rm bulk})^{1/\sigma})$, rounded upward to
a code scale and with an enlarged constant, controls the bulk and
controller-preparation terms. The raw-interface term remains explicit;
fixed positive intensity does not make arbitrary-size once-only raw
interfaces uniformly reliable. For a fixed interface, choose the initial
intensity small enough to put this term inside a fixed total budget, such
as $1/4$.

A kernel may be a universal interpreter step with its actual program,
head/address, work and halt records carried as protected mutable data.
The immutable layout then implements the interpreter and the correction
machinery, rather than a supplied input-dependent answer. A known finite
step bound is counted down; a protected halt flag may suppress later work.
The theorem specifies a channel for every supplied input, including
inputs entangled with an inaccessible reference.

\begin{theorem}[Optimal asymptotic autonomous memory diameter]\label{thm:diam}
Under the same premises, with declared input/output access error strictly
below $1/4$, and at one sufficiently small fixed positive intensity,
there is a uniform family retaining a once-supplied unknown
qubit for a requested time $T$, then delivering it locally with half-diamond
error at most $1/4$. Its post-input duration is $T(1+o(1))$ and
\[
 \diam(T)\leq T^{1/r+o(1)},\qquad
 N(T)\leq T^{2/r+o(1)},\qquad
 U_{\rm all}(T)\leq T^{1+2/r+o(1)}.
\]
Conversely every admissible finite device contained in a box of radius
$L(T)$ and delivering a qubit with this error at time at least $T$ satisfies
$L(T)=\Omega((\lambda T)^{1/r})$.
Thus the logarithmic growth exponent of the full autonomous device's
diameter is exactly $1/(\alpha-d)$. The hardware and work exponents are
achievable upper bounds, not optimality statements.
\end{theorem}

For a Bell input, the delivered state has Bell overlap at least $3/4$,
which exceeds the separable bound $1/2$. This is a quantum input-preserving
operation, not only classical one-bit storage or offline decoding.

\section{Accepted interfaces used, and the new bridges}

We reuse the accepted geometric source theorem and actual noisy ascent
and descent. Projection onto a moving wire gives a Pareto exponent
$a=\alpha-d+1=r+1$ and changes the intensity by a fixed geometric factor.
We absorb that factor into displayed constants. A code block of width
$B^m$ has protected rectangle duration $A^m$ and bad-geometry probability
\begin{equation}\label{eq:pm}
 p_m\leq K_*\lambda\theta^m
       =K_*\lambda H^{\kappa-r},
 \qquad \theta=AB^{-r}<1.
\end{equation}
The separated-source proof gives
$p_m\leq S_B^2p_{m-1}^2+c_B\lambda\theta^m$; it does not assert
independent repair cones or independent effective faults. Geometry is the
same before and after the controller and I/O additions, up to constant
local tracks, padding and timing.

The quantum compiler uses strong nonoverlapping rectangles. Its
deterministic properties are:
\begin{description}
\item[Correctness.] Good rectangles preserve the intended decoded action
on valid inputs, with an arbitrary external reference.
\item[Validity.] A good identity/recovery rectangle maps \emph{any physical
input} into the recursively valid class.
\item[Semi-correctness.] On arbitrary input the decoded gate part has the
intended action after the leading recovery, retaining syndrome and
environment information in the proof.
\end{description}
These are the Rec-Cor, Rec-Val and Rec-SemiCor properties proved recursively
in AGP, Section 10, Lemma 6 \cite{AGP}. The decoder sweep retains syndrome;
it does not assign independent hypothetical logical values to malformed
child states. The proof of Rec-Val applies the base property regardless
of the state emerging from lower-level decoders.

For a fixed child-fault budget $D$, choose spread $s_0$ and correction
capacity at least $2s_0D$. Leading recoveries remove old valid-input errors,
and the new residual is at most $s_0D$ when the total child budget is at
most $D$. The same deterministic induction gives the three properties at
every level. Measurement-free polynomial-code gadgets with arbitrary
fixed correction capacity supply the required base interface; see
ABO \cite{ABO} and the explicit formulation in DST, Lemma 4 and
Definition 4 \cite{DST}. Local routing changes fixed constants only.

In particular, write $\Norm_m$ for a protected identity rectangle. Then
\[
 \Pr(\Norm_m\text{ bad})\leq p_m
\]
even when its input is an arbitrarily damaged physical block. There is
no additional $\sum_{j<m}p_j$ interface charge for this normalization.
This distinction is the central reason reuse can work.

The new work consists of the noisy product-controller initialization,
the mixed control interface, the sacrificial-lane loop, and the
fixed-intensity change of padding used to optimize whole-device diameter.
We give those bridges explicitly.

\section{Noisy initialization and mixed classical/quantum control}

\begin{lemma}[Known product words]\label{lem:init}
In a fixed-duration preparation window, apply noisy local
$\operatorname{RESET}_c$ to the canonical leaves of a hierarchical
classical repetition encoding of a known $c\in\{0,1\}$. At extent $H=B^m$,
the probability that the resulting block violates its recursive validity
promise is at most $C\lambda H^{-r}$.
\end{lemma}

\begin{proof}
This is a spatial source recurrence, not a fresh unknown-input encoding.
Let $q_k$ be the sufficient geometric bad event at level $k$ after the
parallel reset. Truncate actual physical marks at
$R=\rho B^{k-1}$, for the fixed guard constant $\rho$. The repetition
margin is chosen so that parent invalidity requires enough bad canonical
children to include two with disjoint truncated spatial source sets.
Their bad probabilities are each at most $q_{k-1}$. The preparation
duration is constant, so the mean number of large marks touching the
parent wire is at most $c_B\lambda B^{-rk}$. Consequently
\[
 q_k\leq S_{\rm init}^2q_{k-1}^2+c_B\lambda\beta^k,
 \qquad \beta=B^{-r}<1,\qquad q_0\leq c_0\lambda.
\]
Induction at sufficiently small intensity gives
$q_k\leq K_{\rm init}\lambda\beta^k$.
The event depends on mark geometry and holds for arbitrary fault values;
independence is used only for disjoint sets of Poisson sources.
No factor $A^k$ occurs because preparation lasts a fixed physical time,
rather than a level-$k$ logical operation.
\end{proof}

Preparing $w_{\rm c}$ known counter and flag bits has error at most
$C\lambda w_{\rm c}H^{-r}$ by a union bound. It requires at most
$O(w_{\rm c}H)$ local writes, bounded depth under the stated parallel
access, and the actual preparation faults above. Only the designated
cohort's controller leaves need these initial values; other cohorts may
be arbitrary.

This operation is deliberately different from loading $\log K$ raw
unknown bits through the ascent. The latter would add an unwanted
$O(\lambda\log K)$ interface term. The known word $K$ is instead written
as a product repetition pattern, under the same faults, with the spatial
bound just proved.

\begin{lemma}[Mixed control interface]\label{lem:hybrid}
A valid repetition-coded classical bit can control a quantum encoded
gate with a fixed fault budget, using only noisy local gates, private
scratch and a constant-size template. The interface preserves all
classical--quantum correlations and permits $A=CB$ padding with a fixed
finite primitive alphabet.
\end{lemma}

\begin{proof}
One explicit construction first prepares a fault-tolerant encoded
quantum zero. For each quantum-code coordinate separately, reset private
classical vote scratch, read all canonical controller lineages through
dephase-and-copy operations, sort/majority those copies, and use the
result only for the corresponding coordinate's $X^c$. Polynomial-code
logical $X$ is transversal, so correct coordinate votes prepare the
quantum encoding of $c$.

Use classical repetition length $n_{\rm c}=8D+1$ with initial allowance
$2D$. A total of $D$ faulty child locations can contact at most $2D$
additional canonical lineages. Every untouched private vote therefore
still sees at least $n_{\rm c}-4D=4D+1$ correct copies. A damaged vote or
direct data fault affects only a bounded number of quantum coordinates.
No shared mutable majority result controls the whole quantum word.
Canonical sources only undergo copying and routed swaps until the votes
are complete. Finish with an ordinary classical refresh and quantum
recovery. The same total budget covers early and late faults; its
correction margins are enlarged by the fixed spread constant.

The resulting quantum control is a valid encoding of $c$ when the
classical input is valid. Arbitrary incoming scratch is reset. On invalid
controller inputs, final strong quantum recovery supplies a valid but
unspecified quantum block; no intended classical value is asserted.
Use this control in an ordinary protected controlled gate.

All classical reads explicitly dephase. Thus the promise is for classical
or classical--quantum control states, not coherent superpositions of
classical repetition words. Coherent program/control registers can stay
in the quantum code.

The number of private vote rows, code coordinates and routed interactions
is constant. Typed ports specify the classical or quantum code format
in the immutable finite template. A typed swap routes the fixed number
of live child tokens into the other block's slots, carrying their code
format with them, then refreshes/recovers. Recursive routing bottoms out
in physical swaps. One bad child operation contacts only boundedly many
token lineages. Each constant template has duration $O(B)$, so the mixed
compiler retains $A=CB$ with $C$ independent of padding $B$.
\end{proof}

The conditional swap needed below has an exact finite implementation.
For qubits it is a finite reversible Fredkin circuit; for fixed
prime-dimensional registers, finite-field reversible arithmetic supplies
the controlled additions and swaps. Generic approximate universality is
not used with a fixed error on every lap. Algorithm-specific
approximations, if present, are budgeted over the full computation.

\section{The decisive reuse construction}

Each temporal cohort contains a protected data/work register $D$, its
protected classical controller, and a sacrificial I/O register $I$.
The controller contains $\mathrm{loaded}$, $\mathrm{done}$, and a binary
countdown. Initially its known values are $0,0,K$, prepared by
Lemma~\ref{lem:init}. The raw unknown input starts in the compact local
port of $I$. The quantum state of $D$ can be arbitrary.

The same finite superperiod is spatially compiled once and repeated.
It performs the following operations in order.
\begin{enumerate}
\item Run the actual multiscale raw-to-level-$m$ ascent on $I$. During
this stage, $D$ and the controller undergo complete protected identity
rectangles.
\item Apply $\Norm_m$ before coupling $I$ to the retained registers.
Normalize initially arbitrary $D$ as well. If
$\mathrm{loaded}=0$, swap the input from $I$ into the designated $D$
registers, initialize required blank work registers by protected
preparations, and set $\mathrm{loaded}=1$.
\item Execute one protected kernel step, including its mutable program,
head/address and work records. Compute the release flag from
$\mathrm{done}=0$ and $\mathrm{count}=1$, then decrement the positive
count. A done flag suppresses further intended work.
\item Normalize $I$ again before coupling. If release is true, swap the
designated output registers from $D$ into $I$ and set $\mathrm{done}=1$.
Otherwise this coupler acts as the identity on $D$.
\item Run the actual multiscale descent on $I$, ending at the same raw
port at the period boundary. $D$ and the controller continue protected
identities. Read the port only on the specified $K$-th boundary.
\end{enumerate}

Only the first ascent and last descent must succeed. A known-zero reset
of $I$ between its uses is optional and unnecessary. Before release,
conditional swap is the identity on $D$; at release, swap transfers $D$'s
state irrespective of the previous valid logical contents of $I$.
Normalization and its own reset ancillas are nevertheless charged.

\subsection{Why arbitrary intermediate interface damage is harmless}

The key load channel, after normalization, is
\[
 F_{\rm load}
 =|0\rangle\!\langle0|_{\rm loaded}\otimes\SWAP_{ID}
  +|1\rangle\!\langle1|_{\rm loaded}\otimes I_{ID}.
\]
On its identity branch, every joint state $\rho_{IDR}$ satisfies
\[
 \Tr_I\!\left[
 F_{\rm load}(|1\rangle\!\langle1|\otimes\rho_{IDR})
 F_{\rm load}^{\dagger}\right]
 =|1\rangle\!\langle1|\otimes\rho_{DR}.
\]
A trace-preserving normalization on $I$ and its scratch also leaves the
$DR$ marginal unchanged, even when $I$ is correlated with $D$, $R$ and
the environment. Good gate correctness then applies to the valid encoded
ports. At the first load, the accepted successful ascent supplies the
intended unknown state; arbitrary-input normalization alone would not
earn that input access. On release, a swap gives
$\rho^{\rm after}_{IR}=\rho^{\rm before}_{DR}$ regardless of correlations
with the destination.

Between these couplers, arbitrary I/O-lane histories are local channels
on that lane and its environment. Any actual burst that also touches a
data/controller operation is already included in that operation's
geometric fault count. Earlier interface damage does not become an
unmarked later data fault: the entire malformed port is normalized
before the next coupling.

\begin{proof}[Proof of the channel statement in Theorem~\ref{thm:reuse}]
Condition on a geometry for which the controller initialization is valid,
the first ascent succeeds, every counted top-level protected operation
is good, and the last descent succeeds. Intermediate ascents and descents
are unrestricted.

First normalize the arbitrary destination $D$ and use the successful
input ascent to perform the first load. Inductively, after lap $j$,
the retained logical state has undergone exactly $j$ intended kernel
steps, the loaded flag is one, and the counter and done flag have their
ideal values. On later loads the channel just displayed is the identity
on $D$, regardless of the sacrificial state. Good kernel and controller
rectangles perform the next step. Before the last lap the release channel
is also the identity. On the last lap it transfers the desired state to
$I$, and the successful final descent delivers it locally. All identities
hold jointly with a reference and the retained syndrome/environment.

There are $O(KG)$ counted top rectangles, including the normalizers,
couplers, controller computation, blank preparations, and protected
identities during I/O work. Their union failure probability is at most
$O(KGp_m)$. Initialization contributes
$O(\lambda w_{\rm c}H^{-r})$; the first and last interfaces contribute
$C_{\rm io}\lambda$. A mixture over unfavorable geometry adds at most
its probability to half-diamond error, giving
Equation~\eqref{eq:generalerror}. This uses no independence between bad
top rectangles and imposes no reliability condition on the intermediate
raw I/O passes.
\end{proof}

\subsection{Timing and locality of the repeated law}

The ascent and descent can be made exact integer multiples of the
top-level rectangle duration. If the one-level lift has child-depth
$e_{\rm up}$, choose a fixed $c_{\rm up}$ with
$c_{\rm up}(A-1)\geq e_{\rm up}$ and construct
\[
 U_m=U_{m-1}
  +[c_{\rm up}(A-1)-e_{\rm up}]
       \text{ protected child periods}
  +\text{encoder},
 \qquad |U_m|=c_{\rm up}A^m.
\]
The extra waits occur after the smaller-scale input is protected.
Similarly perform recovery and inverse encoding first, then the
necessary protected child waits, then the smaller descent, obtaining
$|D_m|=c_{\rm down}A^m$. The final physical access is at the end.
Padding preserves the geometric-sum interface bound and gives whole
top-level identity windows for retained data/controller registers.
Additional annulus-width padding is also placed while the relevant
quantum state remains encoded.

Place one physical superperiod, of $P$ rows and width $W$, on the accepted
two-track conveyor. Every tick applies the identical local channel
\[
 F=G_{\rm spatial}\circ
   \prod_y\SWAP(B_y,A_{y+1})\circ
   \prod_y\SWAP(A_y,B_y).
\]
Indices are modulo $P$ and $G_{\rm spatial}$ is the fixed row-gate array.
At complete ticks,
\[
 A_y(t+1)=G_y(A_{y-1}(t)),\qquad
 B_y(t+1)=B_{y+1}(t).
\]
The desired I/O, data and controller therefore remain in one temporal
cohort. Other cohorts never ideally couple to it. The identity is exact
for arbitrary and even correlated junk; fault tubes include both
operands during the two swap microphases.

Pad to $P\geq cW$ and use the annular embedding
\[
 (x,y)\longmapsto
 (-R_0+(R_0+x)\cos(2\pi y/P),\,
        (R_0+x)\sin(2\pi y/P)),\qquad R_0=c_1P.
\]
It has bounded local edge length and bounded register density. At row
zero the coordinates are $(x,0)$, so the compact raw I/O port is independent
of $P$ and $H$. The same row-zero port is reached each lap. Constant
interleaved register tracks implement multiple fixed input/output cells
without horizon-dependent raw-port separation.

At each time a rigid transformation straightens the active radial row.
The moving interaction tube has bounded thickness because $P\geq cW$.
Transforming the physical Poisson process at each time preserves its
space-time intensity; dominating cubes by fixed-factor balls and
projecting gives the same effective wire-tail exponent $r+1$. Repeatedly
visiting the same physical positions creates no independence assumption:
the accepted source recurrence bounds each rectangle, and the proof
uses a union bound over successive uses.

\section{Execution costs and the memory specialization}

The compiled width is $W=O(wH)$. One superperiod contains the fixed
ascent/descent stages, all the kernel/controller work, normalizations,
conditional couplers and protected padding. Its duration is
$P=O(\bar bH^\kappa)$ with $\bar b=O(b+w+1)$.
One annulus therefore has $N=O(WP)$ registers, including both transport
tracks and all off-cohort rows. Running it $K$ times takes $KP$ ticks and
$O(NKP)$ register updates, including idle registers. These are
Equations~\eqref{eq:res1}--\eqref{eq:res2}.

For the same width, code level and $K$-period task, a complete spatial
unroll would use $O(KWP)$ registers and $O(K^2WP^2)$ all-register updates.
Reuse removes a factor $K$ from both, with the same post-input duration.
This comparison keeps the encoding scale and controller costs equal.

Uniform layout generation takes $O(N\operatorname{polylog}N)$ work
using the recursive templates, with constants of the selected fixed
family included. Fabrication installs $O(N)$ bounded-size components
and local connections; serial fabrication takes $O(N)$ time. The
theorem does not supply autonomous fabrication. Controller-product
preparation costs $O(w_{\rm c}H)$ local reset operations in bounded
parallel depth. Fresh ancillas, their reset/disposal operations, and all
subsequent noisy work are contained in the compiled execution count.
Final delivery is at a fixed finite raw port; its constant number of
access/readout operations and their errors remain explicit. These are
abstract channel-operation costs, not measured power, cooling cost, or
experimental timing guarantees.

For memory, take the kernel to act as the identity on the unknown qubit.
Choose the counter width in advance as
$w(T)=\lceil\log_2(T+2)\rceil+O(1)$, so it can represent the eventual
number of laps. A nearest-neighbor ripple countdown and its tests have
$b=O(w)$, hence $\bar b=O(w)$ and $G=O(w^2)$.
Given the physical superperiod $P$, set $K=\lceil T/P\rceil$. Then
\[
 T\leq KP<T+P.
\]
The counter circuit's width and schedule do not depend on its initialized
value $K$, so this choice is not circular.

Using $p_m=O(\lambda H^{\kappa-r})$,
\begin{align}
 KGp_m
 &\leq C\lambda\left(\frac{T}{wH^\kappa}+1\right)
                    w^2H^{\kappa-r}\\
 &\leq C\lambda\bigl(TwH^{-r}+w^2H^{\kappa-r}\bigr).
\end{align}
The first term is the accumulated exposure; the second covers integer
rounding and one whole extra superperiod. Both are controlled by taking
\[
 H\ \geq\
 \max\left\{H_{\min},
 \left(\frac{C\lambda Tw}{\epsilon_{\rm bulk}}\right)^{1/r},
 \left(\frac{C\lambda w^2}{\epsilon_{\rm bulk}}\right)^{1/(r-\kappa)}
 \right\}
\]
and rounding upward to the next physical code scale. For each fixed
family the first growing term dominates eventually, so
$H=\widetilde O(T^{1/r})$. Initialization is smaller:
$O(\lambda wH^{-r})$.

Therefore each fixed padding family gives
\begin{align}
 W&=\widetilde O(T^{1/r}),&
 P&=\widetilde O(T^{\kappa/r}),&
 \diam&=\widetilde O(T^{\kappa/r}),\\
 N&=\widetilde O(T^{(1+\kappa)/r}),&
 U_{\rm all}&=\widetilde O(T^{1+(1+\kappa)/r}).
\end{align}
Since $\kappa<r$, its period is $o(T)$ and post-input duration is
$T(1+o(1))$. The constants and polylogarithms depend on the fixed
compiler. Construction latency is separately charged and can exceed
the subsequent retention time.

\section{Changing padding at one fixed positive intensity}

Simply taking $B\to\infty$ in a threshold statement would be insufficient:
its allowable physical intensity might shrink. We instead use baseline
levels to enter each new padding scheme, keeping the physical burst law
and intensity fixed.

Fix the codes, local dimensions, mixed primitive alphabet and fault
budget once. Choose a baseline padding $B_0$ with $A_0=CB_0$ and
\[
 \theta_0=A_0B_0^{-r}<1,\qquad
 p^{(0)}_\ell\leq K_0\lambda\theta_0^\ell.
\]
Choose the physical $\lambda$ below this one baseline threshold and
below the fixed raw-interface budget. This choice precedes all requested
horizons and all later exponent refinements.

For a selected upper padding $B$ with
$A=CB$ and $\theta_B=AB^{-r}<1$, use $\ell_0$ baseline levels first:
\[
 L_0=B_0^{\ell_0},\qquad \tau_0=A_0^{\ell_0},\qquad
 \mu=\lambda\tau_0L_0^{-r}=\lambda\theta_0^{\ell_0}.
\]
The base child failure satisfies $p_0\leq K_0\mu$. The code and all
semantic gadget contracts are identical across this seam; only spatial
padding changes. No logical code conversion or supplied encoded input
is required. Extra baseline levels are executed by the actual noisy
ascent.

\begin{lemma}[Source recurrence across the seam]\label{lem:seam}
At upper level $j$, with width $L_j=L_0B^j$ and duration
$\tau_j=\tau_0A^j$, the actual physical Poisson source process gives
\[
 p_j\leq S_B^2p_{j-1}^2+c_B\mu\theta_B^j.
\]
For sufficiently large finite $\ell_0$, it follows that
$p_j\leq K_B\mu\theta_B^j$ at the original physical intensity.
\end{lemma}

\begin{proof}
Truncate physical marks at radius $\rho L_0B^{j-1}$. If the parent is
bad without a large mark, the accepted separated-source combinatorial
lemma supplies two bad child rectangles with disjoint truncated source
sets. Restricting a baseline child's source set only decreases its
geometric badness, so its probability is bounded by $p_0$ at the seam.
Independence follows from the two disjoint sets of actual Poisson marks.
The aspect ratio $\tau_0/L_0$ causes no extra dependency: spatial guards
and temporal windows scale with their respective lengths, and marks
are instantaneous. The large-mark mean is bounded by
\[
 c_B\lambda\tau_jL_j^{-r}
 =c_B\mu(AB^{-r})^j.
\]
These two alternatives give the recurrence.

Take $K_B\geq\max(K_0,2c_B)$ and
$\mu\leq\theta_B/(2S_B^2K_B)$. Substituting the induction hypothesis,
the quadratic term divided by $K_B\mu\theta_B^j$ is at most
$S_B^2K_B\mu\theta_B^{j-2}\leq1/2$ for $j\geq1$.
The tail term is also at most $1/2$ of the target. Since
$\mu\to0$ as $\ell_0\to\infty$, a finite boost supplies this condition.
\end{proof}

Here $\mu$ is a tail prefactor and a bound on baseline child failures.
It is \emph{not} the intensity of a newly independent Poisson process of
coarse faults. Numerous small physical marks can occur inside a baseline
block. Their aggregate effect is handled by the baseline source theorem.

The interface budgets can also be made uniform across upper paddings.
The baseline portion of first ascent and last descent costs at most
\[
 C_{\rm base}\sum_{\ell=0}^{\ell_0}p^{(0)}_\ell
 \leq \frac{C_{\rm base}K_0\lambda}{1-\theta_0},
\]
independently of $\ell_0$. If $M_B$ bounds the number of child operations
used by an upper lift/descent stage, its whole suffix costs at most
\[
 \frac{M_BK_B\mu}{1-\theta_B}.
\]
Choose $\ell_0$ large enough that this is below one fixed
$\epsilon_{\rm suffix}$, as well as satisfying the recurrence condition.
This changes the layout and executed ascent but neither the initial
unknown seed nor the physical intensity.

The initial known controller words have the spatial-only prefactor
$\lambda L_0^{-r}$. The same spatial recurrence gives
$q_m\leq C_B\lambda(L_0B^m)^{-r}$; increase the finite boost if needed
for its threshold. The duration prefactor $\tau_0$ does not appear in
this product-preparation bound.

Writing $H=L_0B^m$, upper failure is
$p_m\leq C_B\lambda\tau_0L_0^{-r}A^mB^{-rm}$ and one period is
$P=O_B(\bar b\tau_0A^m)$. The accumulated memory term is therefore
still $O_B(\lambda T\,\operatorname{polylog}T\,H^{-r})$:
the time dilation in exposure cancels against the number of periods.
The preceding fixed-family resource exponents remain unchanged.

\subsection{A uniform algorithm attaining the limiting exponent}

Choose effectively increasing $B_j$ so that
$\kappa_j=\log_{B_j}(CB_j)\downarrow1$ with $\kappa_j<r$.
For each $j$, compute a finite baseline depth satisfying
Lemma~\ref{lem:seam} and the common suffix budget. Each complete fixed
family has finite computable constants in its resource and error bounds.

Choose increasing thresholds $T_j$ such that, for all $T\geq T_j$,
its constants and logarithmic factors are bounded by $T^{\varepsilon_j}$,
where $\varepsilon_j\downarrow0$, its baseline dimensions fit those
bounds, and $P_j(T)\leq T/j$. On $[T_j,T_{j+1})$ use that complete
$j$-th family. Then
\begin{align}
 \diam(T)&\leq T^{\kappa_{j(T)}/r+\varepsilon_{j(T)}}
                    =T^{1/r+o(1)},\\
 N(T)&\leq T^{(1+\kappa_{j(T)})/r+\varepsilon_{j(T)}}
                    =T^{2/r+o(1)},\\
 U_{\rm all}(T)&\leq T^{1+(1+\kappa_{j(T)})/r+\varepsilon_{j(T)}}
                    =T^{1+2/r+o(1)}.
\end{align}
The period condition gives $KP=T(1+o(1))$ uniformly through the
diagonal. All error budgets are the same across families.

There is no need to hide an unbounded search in layout selection.
Enumerate certified family records for at most $O(\log T)$ computation
steps and use the largest completed eligible record. Every fixed record
eventually completes and becomes eligible, so the selected index still
tends to infinity. Conservative thresholds also absorb family-dependent
layout-generation constants. The finite physical gate alphabet remains
unchanged.

For rational implementation bookkeeping one can instead use certified
$1<r_j<r$ increasing to $r$ and design each family for the heavier tail
$d+r_j$. The baseline may use the first heavier design law for a uniform
small-error bound. At each seam use the maximum of the baseline failure
prefactor and the upper large-mark prefactor; both tend to zero as the
boost depth grows. Their ratio is a fixed family constant. The same
argument then gives exponents $\kappa_j/r_j$,
$(1+\kappa_j)/r_j$, and $1+(1+\kappa_j)/r_j$, with the displayed limits.
This permits certified finite arithmetic without assuming free exact
real-number computation.

\section{The matching full-device obstruction}

Suppose every mutable register and wire of a device lies in a fixed
box of sup-norm radius $L\geq1$. Marks completely covering this box
arrive at rate
\begin{align}
 q_L
 &=\lambda 2^d\mathbb E(R-L)_+^d\\
 &=\lambda 2^d
   \frac{\Gamma(d+1)\Gamma(\alpha-d)}{\Gamma(\alpha)}
   L^{d-\alpha}
 =c_{\alpha,d}\lambda L^{-r}.
\end{align}
Choose an admissible fault assignment that resets all mutable registers
on every such mark and supplies no input-dependent information afterward.
On the event of a complete reset, which has probability
$p=1-e^{-q_LT}$, the output channel is a replacement channel with
respect to the once-supplied input. Immutable input-independent laws,
layouts and controller initialization cannot reconstruct the erased
unknown state.

Test the channel on one half of a Bell pair. Its Bell overlap is at most
$1/4$ on this erasure branch and at most one on the complementary branch.
Consequently
\[
 \frac12\|\mathcal E_T-\operatorname{id}\|_\diamond
 \geq \frac34(1-e^{-q_LT}).
\]
Error at most $1/4$ requires $q_LT\leq\log(3/2)$, hence
\[
 L\geq
 \left(\frac{c_{\alpha,d}\lambda T}{\log(3/2)}\right)^{1/r}.
\]
Combining this lower bound with the diagonal upper bound proves
Theorem~\ref{thm:diam}. It is a lower bound on the entire device extent,
not merely a selected active wire. It does not imply a matching lower
bound on installed register count for arbitrary sparse layouts.

The accepted complete-cone obstruction for $\alpha\leq d+1$ also
continues to apply: initial controller words and static geometry are
input-independent, and the local update has a uniform finite signal
speed. We do not reprove that accepted boundary here.

\section{Existing-method comparison and limits of the conclusion}

The closest autonomous comparator located is
D\"unnweber--Styliaris--Trivedi (DST) \cite{DST}: a spatially
translation-invariant quantum automaton with noisy control/address
records, under a local fault-path noise condition and a suitably encoded
initial state. Its discussion identifies noisy preparation of the
hierarchy as open; the continuous-time threshold extension is outlined.
Its simulator has $T=T_0T_{\rm sim}=O(M^2)$ for linear block size $M$.

These facts suggest an alternative but do not establish the present
Pareto operation. In a direct large-mark transfer at level $\ell$, the
exposure would contain
\[
 T^\ell M^{-\ell r}=(TM^{-r})^\ell.
\]
With the displayed dilation this asks for roughly $r>2$. This is an
obstruction to that \emph{direct transfer argument}, not a no-go theorem
for all versions of the automaton. Moreover a control-damage proof needs
its actual shrinking/merging repair process dominated by the physical
source geometry. Workload tightness or an offline surviving path does
not supply that map. Our repeated patterned conveyor avoids this
unearned transfer while retaining the strict regime $r>1$.

Classical Toom/G\'acs fault tolerance and ordinary quantum threshold
theorems remain powerful machinery at their hypotheses. The new result
here uses their deterministic semantic principles where established,
then applies the accepted separated \emph{physical-source} recurrence.
It does not replace correlated damage by an iid effective error model.

The precise scope is as follows.
\begin{itemize}
\item The result is a mathematical local-channel construction. It retains
immutable patterned gate laws and synchronous elementary updates.
Translation invariance, noisy asynchronous timing, an autonomous
Hamiltonian, and natural thermal implementation are not established.
\item Initialization includes a finite noisy product-write access to
known controller leaves. Quantum work and sacrificial registers may be
arbitrary. Completely arbitrary controller-state self-initialization is
outside the theorem.
\item The useful endpoint is a fixed finite local quantum/classical
interface, at a known final tick, with a fixed nontrivial total error.
Raw access errors and initial acquisition are explicit. No claim of
arbitrarily small total error at a fixed noisy output cell is made.
\item All hardware is finite for each horizon. The initial
input-dependent seed is fixed, while controller size and layout may
depend on the horizon.
\item The exact diameter \emph{exponent} uses a computable slow choice of
padding. Its $T^{o(1)}$ overhead and constants may be very large.
It is not a practical threshold estimate or an exact
$O(T^{1/(\alpha-d)})$ construction.
\item Installed hardware and update exponents are upper bounds. The
complete-reset lower bound proves only diameter optimality.
\end{itemize}

\section{Executable finite checks}

The accompanying \texttt{reusable\_processor.py} checks the new adapter's
finite channel and transport identities. It does not implement the
physical concatenated code or estimate a noise threshold.

Eight transport cases use periods $7,19$ and lap counts $1,2,5,17$.
They track every temporal cohort on both tracks, verify that I/O,
data and controller registers in each row share one cohort, check every
cohort's return after each lap, and compare the repeated row-gate history
against a complete $K$-period unroll.

Sixteen quantum cases use $K=1,2,5,17$, mutable programs selecting $H$
or $TH$, and two initial-state families: a Bell input with entangled
junk, and a fully correlated four-qubit state. The register order is
I/O, retained data, external reference, environment. Arbitrary channels
on I/O plus environment model failed intermediate interfaces. Each
case has both reset-containing and no-reset variants. The final
I/O--reference marginal is compared with the intended $K$-step channel.
Additional laps test the one-shot guards; output fidelity is claimed
only at the specified delivery phase.

\begin{center}
\begin{tabular}{lr}
\toprule
Check & Maximum numerical error\\
\midrule
Delivered I/O--reference trace distance & $4.853\,10^{-15}$\\
No-reset variant delivery & $1.896\,10^{-15}$\\
Additional quiet-port laps & $4.565\,10^{-16}$\\
Data--reference change under port-only channels & $5.551\,10^{-16}$\\
Density-matrix numerical validity & $1.555\,10^{-14}$\\
\bottomrule
\end{tabular}
\end{center}

For the largest symbolic transport example, period $19$, $K=17$ and
six register columns, reuse needs 228 registers and 73,644 full-rule
register updates. Full unrolling needs 3,876 registers and 1,251,948
such updates. Both take 323 full-rule ticks. Primitive microphase
updates, including idle locations, are separately counted in the JSON.
Physical counter words and code blocks are expanded in the analytical
resource theorem, not represented by one physical register in this
small symbolic check.

\section{Resulting consequence}

The constructive advance is a reusable input-preserving processor whose
noisy controller governs admission, computation count and final local
release. Its input/output lane may fail on every irrelevant pass.
The strong semantic interface makes those failures disappear from the
useful operation's error budget.

For storage, one period is therefore enough hardware. At a single fixed
physical intensity, the padding seam and slow uniform choice of scales
reduce the complete device's diameter to the exponent forced by a
whole-device information erasure. In particular, when $d=2$ and
$\alpha=4$, the achieved bounds are diameter $T^{1/2+o(1)}$, hardware
$T^{1+o(1)}$, and all-register updates $T^{2+o(1)}$, with a fixed
finite unknown input and fixed nontrivial quantum output fidelity.
These are asymptotic theorem bounds; all stated construction and
initialization costs remain paid.

\begin{thebibliography}{9}
\bibitem{AGP}
P. Aliferis, D. Gottesman, J. Preskill,
\emph{Quantum accuracy threshold for concatenated distance-3 codes},
especially Section 10, Lemma 6, pp.\ 43--45.
\url{https://arxiv.org/pdf/quant-ph/0504218}.

\bibitem{ABO}
D. Aharonov, M. Ben-Or,
\emph{Fault-Tolerant Quantum Computation With Constant Error Rate},
extended version; polynomial codes and measurement-free local gadgets.
\url{https://arxiv.org/pdf/quant-ph/9906129}.

\bibitem{DST}
G. D\"unnweber, G. Styliaris, R. Trivedi,
\emph{Quantum Memory and Autonomous Computation in Two Dimensions},
arXiv:2601.20818v3, 25 September 2026; Definition 4,
Lemmas 4 and 9, Sections V--VII.
\url{https://arxiv.org/pdf/2601.20818}.

\bibitem{Accepted}
Accepted internal programme machinery, as supplied in the governing
conversation: strict-Pareto finite-seed scheduled compiler with forced
outputs and separated Poisson sources; noisy ascent/descent and endpoint
extent theorem; reference-inclusive quantum extension and one-pass
two-track autonomous conveyor. This report reuses those stated scopes
and proves the additional compatibility bridges above.
\end{thebibliography}
\end{document}
